
\section{Truncation technique}
\label{sec:truncation}

We introduce the SVD of a Grassmann tensor, 
which is equivalent to the SVD for the corresponding coefficient tensor. 
Let $\mathcal{T}_{\Psi\Phi}$ be a Grassmann tensor whose rank is $2N$. 
We represent the coefficient tensor of $\mathcal{T}_{\Psi\Phi}$ as $2^{N}\times2^{N}$ 
matrix $T_{IJ}$ with $I=(i_{1},\cdots,i_{N})$ and $J=(i_{N+1},\cdots,i_{2N})$. 
Since the Grassmann parity of $\mathcal{T}_{\psi\phi}$ is even, $T_{IJ}$ takes a non-zero value if and only if
\begin{align}
	\sum_{k=1}^{2N}i_{k}\mod2=0
\end{align}
is satisfied. This condition allows us to obtain a block diagonal matrix representation for $T_{IJ}$. According to Ref.~\cite{Kadoh}, we now define the following decimal numeral system,
\begin{align}
	I=
	\begin{cases}
		\sum_{k=1}^{2N}2^{k-1}i_{k}~~~&(i_{2}+\cdots+i_{2N}\mod2=0) \\
		1-i_{1}+\sum_{k=2}^{2N}2^{k-1}i_{k}~~~&(i_{2}+\cdots+i_{2N}\mod2=1).
	\end{cases}
\end{align}
Thanks to this decimal numeral system, the parity of $\sum_{k=1}^{2N}i_{k}$ corresponds with that of $I$. Applying this system for $I$ and $J$ in $T_{IJ}$, one obtains the block diagonal matrix
\begin{align}
	T=
	\begin{bmatrix}
		T^{\mathrm{E}} & 0 \\
		0 & T^{\mathrm{O}}
	\end{bmatrix}
	.
\end{align} 
The SVD for $T$ is obtained from that for $T^{\mathrm{E}}$ and $T^{\mathrm{O}}$,
\begin{align}
\label{eq:even}
	T^{\mathrm{E}}_{IJ}=\sum_{K:\mathrm{even}}U^{\mathrm{E}}_{IK}\sigma^{\mathrm{E}}_{K}V^{\mathrm{E}}_{JK},
\end{align}
\begin{align}
\label{eq:odd}
	T^{\mathrm{O}}_{IJ}=\sum_{K:\mathrm{odd}}U^{\mathrm{O}}_{IK}\sigma^{\mathrm{O}}_{K}V^{\mathrm{O}}_{JK}.
\end{align}
Picking up the largest $D$ numbers of singular values and corresponding singular vectors, $T$ is approximated with a lower-rank matrix. 

We apply the above technique for $MM^{\dag}$, where $M$ is the coefficient tensor in Eq.~\eqref{eq:M_HOTRG}. In Eq.~\eqref{eq:isometry}, the Grassmann isometry defines new bond Grassmann numbers $\Phi$ and $\bar{\Phi}$. It is worth emphasizing that the parity of these Grassmann numbers corresponds to the parity of $K$ in Eqs.~\eqref{eq:even} and \eqref{eq:odd}. This property is significantly useful in developing the current Grassmann HOTRG.













